\documentclass[10pt,a4paper]{article}
\usepackage[margin=2.6cm]{geometry}
\usepackage[T1]{fontenc}
\usepackage{amsmath,amsthm}
\usepackage{newtxtext,newtxmath}
\usepackage{mathpartir,stmaryrd,microtype,booktabs,enumitem,xcolor}
\usepackage{tikz}
\usepackage{needspace}
\usepackage[colorlinks,linkcolor=black,citecolor=black,urlcolor=black]{hyperref}

\newtheorem{theorem}{Theorem}
\newtheorem{proposition}[theorem]{Proposition}
\theoremstyle{definition}
\newtheorem{definition}[theorem]{Definition}
\newtheorem{example}[theorem]{Example}
\newtheorem{assumption}[theorem]{Assumption}
\theoremstyle{remark}
\newtheorem{remark}[theorem]{Remark}

\newcommand{\bx}[1]{\Box_{#1}}
\newcommand{\sem}[1]{\llbracket #1\rrbracket}
\newcommand{\sep}{\mathrel{\#}}
\newcommand{\kw}[1]{\mathsf{#1}}
\newcommand{\keep}{\kw{keep}}
\newcommand{\init}{\kw{init}}
\newcommand{\extend}{\kw{extend}}
\newcommand{\grow}{\kw{grow}}
\newcommand{\any}{\kw{any}}
\newcommand{\val}{\kw{values}}
\newcommand{\st}{\kw{spines}}
\newcommand{\vw}{\kw{views}}
\newcommand{\LV}{\mathit{Loc}_{\mathsf{view}}}
\newcommand{\rigid}{\kw{rigid}}
\newcommand{\fresh}{\kw{fresh}}
\newcommand{\old}{\kw{old}}
\newcommand{\dense}{\kw{dense}}
\newcommand{\apart}{\kw{apart}}
\newcommand{\len}[1]{|#1|}
\newcommand{\tr}[1]{\xrightarrow{#1}}
\newcommand{\resp}{\mathrel{\triangleleft}}
\newcommand{\hoare}[3]{\{#1\}\;#2\;\{#3\}}
\newcommand{\hoarex}[4]{\{#1\}\;#2\;\{#3\mid #4\}}
\newcommand{\code}[1]{\texttt{#1}}
\newcommand{\kF}{\mathsf{field}}\newcommand{\kC}{\mathsf{cell}}\newcommand{\kE}{\mathsf{elem}}\newcommand{\kS}{\mathsf{spine}}\newcommand{\kI}{\mathsf{id}}
\newcommand{\LF}{\mathit{Loc}_{\kF}}\newcommand{\LC}{\mathit{Loc}_{\kC}}\newcommand{\LE}{\mathit{Loc}_{\kE}}\newcommand{\LS}{\mathit{Loc}_{\kS}}\newcommand{\LI}{\mathit{Loc}_{\kI}}
\newcommand{\lean}[1]{\texttt{\small #1}}
\newcommand{\status}[1]{\textsf{\small[#1]}}
\newcommand{\cC}{\mathcal{C}}

\title{Read-only frames\\[2pt]\large Facts of a component across calls of code outside it}
\author{TScert working note}
\date{18 November 2026}

\begin{document}
\maketitle

\begin{abstract}
\noindent
A function that calls a callback of unknown effect keeps, by the rules of
\emph{Tolerant frames}~\cite{tf}, no fact about the heap across the call.
This note states a rule that keeps facts about a component's own storage
across such calls: a component is a set of functions entered only through
its interface, and the callback, which reaches the component's storage
only at read-only types, changes it only by calling the component's
members, whose effects bound it. The rule, \textsc{Read-only frame},
bounds a call of outside code by the members' protocol on the component's
protected region, under an assumption on unanalysed clients, and a
proposition carries a fact across the call by its support. Its checks read
a flow of allocation sites, an escape analysis that an untrusted procedure
proposes and a decision procedure, proved sound, checks. The calculus,
effects, supports, the frame rule and Reliance are those of
\emph{Tolerant frames}; each statement here carries its status mark in
that note's convention: \status{Lean} names the theorem that proves it,
\status{conj.} marks a conjecture, and \status{assumed} a specification
assumed of the host or of clients.
\end{abstract}

\section{The problem}\label{sec:problem}

In the module
\begin{center}\small
\begin{tabular}{@{}l@{}}
\code{type H = \{ readonly widths: readonly number[] \};}\\
\code{export function prepare(n: number): H \{ return \{ widths: new Array<number>(n).fill(1) \}; \}}\\
\code{export function walk(h: H, visit: () => void): number[] \{}\\
\code{~~const w = h.widths, n = w.length, out: number[] = [];}\\
\code{~~for (let i = 0; i < n; i++) \{ visit(); out.push(w[i]!); \}}\\
\code{~~return out; \}}
\end{tabular}
\end{center}
the read \code{w[i]!} needs $i<\len w$ after \code{visit()}, whose
effect $\top$ allows shortening \code{w}. Yet the visitor holds \code{h}
only at the read-only type \code{H}, and changes the module's storage
only by calling \code{prepare} or \code{walk}. This note states that
argument as a rule. It builds on \emph{Tolerant frames}~\cite{tf}: its
heap of locations by kind, changes and their order
$\keep\sqsubseteq\init\sqsubseteq\kw{scalar}\sqsubseteq\extend\sqsubseteq\grow\sqsubseteq\any$,
effects and their action, supports and respect (its Sections~3 and~4),
the rule \textsc{Frame} and Proposition Reliance (its Section~6).

\section{Concepts}\label{sec:concepts}

Six concepts state the argument, each shown on the module above.
\begin{description}[leftmargin=1.5em,style=nextline]
\item[Component] A set $\cC$ of functions entered only through its
\emph{interface}, here the module's exports. Its \emph{members} are the
functions outside code can call: the interface's, and each function of
$\cC$, closure or not, that reaches outside code as a value. A call runs
\emph{outside} $\cC$ when its callee is not in $\cC$, as the invocation's
component (\emph{Tolerant frames}, Section~4) records.
\item[Crossing] A value \emph{crosses} \emph{in} as an entry's argument,
a callback's result or a caught value, and \emph{out} as an entry's
result, a callback's argument, a thrown value or a value stored where
outside code reaches it. A type is \emph{read-only} when its fields are \code{readonly}, its
arrays \code{readonly} arrays and its typed arrays \code{ArrayLike}, and
the types of its positions are read-only in turn; a field a type does not
mention has type \code{unknown}, which is read-only. A location is
\emph{reached read-only} from a crossing at type $T$ when some prefix of
a chain of field and element reads from the value to it ends at a
position of read-only type. Everything at that position and below it is
then reached read-only: from \code{h} at \code{H}, the record's field location
\code{h.widths} and the array's spine and cells. A site that one crossing
reaches writably is released, whatever other crossings do; released or
outside storage that some crossing in reaches read-only is protected,
whatever other crossings do. A position whose type does not bound what
arrives there accepts every value: a callback's result typed
\code{void}, which a function returning any value meets, a caught value,
and a position of type \code{any}; a value that crosses in there is read
at no read-only type.
\item[Marks] Each allocation site of $\cC$'s code, a native call's
included, is \emph{local} when its storage is held only by $\cC$'s
frames (locals and arguments) and by local storage, \emph{internal} when
no crossing value reaches it, \emph{exposed} when crossing values reach
it only read-only, and \emph{released} otherwise. Storage of outside code
is \emph{outside}. A site's mark is proposed for it and confirmed by the
check (K) below, against every crossing and store of $\cC$'s code. \code{prepare}'s array is exposed; \code{out} is
released.
\item[Protected region] $R_\cC$ holds local and exposed storage, and the
released and outside storage that a crossing in reached read-only, from
then until the outermost active call of $\cC$ returns, and afterwards
while $\cC$ \emph{holds} it: while a frame, a module or closure variable,
or storage of $\cC$ references it or storage it is reachable from.
\item[Protocol] A protocol $\Psi$ gives a change for each \emph{part}
$\mathrm{part}(l)$: the spines, the cells, or the fields at one key, of
one site or of outside storage. The spine of \code{prepare}'s array has
the part (the site \code{new Array}, spines), and the field
\code{widths} of its record the part (the record's site, fields at
\code{widths}); all outside storage of one kind and key is one part,
such as every outside spine. $\Psi$ \emph{bounds the members} when
each member, its own calls of outside code bounded by the rule at
$\Psi$, changes each location of $R_\cC$ existing at its entry within
$\Psi$ at its part, at both exits. The members' \emph{guarantee}
$G_\cC$ is a set of invariants each keeps at every write to their
support and at both exits; an invariant is written \code{/** @invariant
\ldots{} */} on a type of the interface, over \code{this}, and holds of
every value of that type. Here $\Psi_\cC=\keep$: both members write
only storage they allocated.
\item[Client obligation] $\Theta_\cC$, a protocol stated by $\cC$'s API,
bounds what clients change of $R_\cC$. A read-only type at the
interface, or no statement, means $\Theta_\cC=\keep$. The protocol
$\Psi_\cC\sqcup\Theta_\cC$ gives each part the join, in the order of
changes, of the two changes at that part.
\end{description}

\begin{assumption}[Clients]\label{ass:clients}
Code outside $\cC$ that is not analysed enters $\cC$ only through its
interface, passes at each position of the interface a value of the
position's declared type, and each of its own steps changes each location
of $R_\cC$ within $\Theta_\cC$ at its part
\status{assumed, of each component's clients. Lean: the hypothesis \lean{OutsideStep.row}; over the calculus's steps, \lean{ClientFrame.Keeps}; entry only through the interface as code provenance, every code tagged with $\cC$ one of $\cC$'s own, \lean{ClientFrame.Provenance}, kept by every step, \lean{step\_prov}; and a unit accepted under it is safe for the executions keeping it, \lean{FramedRequestSafe}}.
\end{assumption}

\noindent
Three examples show the marks and the protected region at work; each is
accepted.

\begin{example}[A local cursor]\label{ex:rocursor}\mbox{}
\begin{center}\small
\begin{tabular}{@{}r@{\hspace{1em}}l@{}}
1 & \code{type Visitor = (w: number) => Modifies<\{ values: "any"; spines: "grow" \}, void>;}\\
2 & \code{export function walk(xs: number[], visit: Visitor): number \{}\\
3 & \code{~~const c = \{ i: 0 \}; let s = 0;}\\
4 & \code{~~while (c.i < xs.length) \{ visit(c.i); s += xs[c.i]!; c.i++; \}}\\
5 & \code{~~return s; \}}\\
\end{tabular}
\end{center}
\emph{Accepted.} The visitor's type states its effect,
$\Phi=\{\val:\any,\ \st:\grow\}$. The cursor's site (line 3) is local: only
\code{walk}'s frame holds \code{c}, and the visitor receives the number
\code{c.i}. Across \code{visit(c.i)} on line 4 the obligation
$c.i<\len{\mathit{xs}}$, with $0\le c.i$ from the loop, has the support
$\{c.i:\keep,\ \mathit{xs}.\kw{spine}:\grow\}$, tested location by
location: $c.i$ is in $R_\cC$, where $\Psi_\cC\sqcup\Theta_\cC=\keep$;
the spine of \code{xs}, which crossed in at the writable type
\code{number[]}, is outside $R_\cC$, where $\Phi$ allows $\grow$.
\end{example}
\begin{example}[A released result]\label{ex:roreleased}\mbox{}
\begin{center}\small
\begin{tabular}{@{}r@{\hspace{1em}}l@{}}
1 & \code{const held: number[][] = [];}\\
2 & \code{export const rows = (): number[][] => held;}\\
3 & \code{export const fill = (n: number): number[] => \{}\\
4 & \code{~~const row: number[] = []; for (let i = 0; i < n; i++) row.push(1);}\\
5 & \code{~~held.push(row); return row; \};}\\
6 & \code{export const count = (n: number): number => held.length + n;}\\
\end{tabular}
\end{center}
\emph{Accepted.} The sites of \code{held} (line 1) and \code{row}
(line 4) are released: \code{rows} and \code{fill} give them out at
writable types. The entries take only numbers, so no released array comes
back in, none is protected, and the store of line 5 is into storage
outside $R_\cC$, which (O) does not constrain.
\end{example}
\begin{example}[Storage passed in read-only]\label{ex:roin}\mbox{}
\begin{center}\small
\begin{tabular}{@{}r@{\hspace{1em}}l@{}}
1 & \code{const mine = [1, 1, 1];}\\
2 & \code{walk(\{ widths: mine \}, () => \{\});}\\
\end{tabular}
\end{center}
\emph{Accepted}, with \code{walk} of Section~\ref{sec:problem}. The
client's array is outside storage. The crossing of the record at
\code{H} on line 2 reaches it read-only, so it is in $R_\cC$ from that
crossing until \code{walk} returns, and $i<\len w$ survives the visit as
in Example~\ref{ex:ro}. By Assumption~\ref{ass:clients} the client is
bound not to change \code{mine} during the call, even through its own
writable alias: a visitor that pops \code{mine} breaks the assumption.
\end{example}

\section{The rule}\label{sec:rule}

\begin{mathpar}
\inferrule*[left=Read-only frame]{c:\Phi\mid\Phi^{e} \\ c\text{ a call, inside a call of }\cC\text{, of code outside }\cC \\\\ \Psi_\cC\text{ bounds the members} \\ \text{(S)} \\ \text{(K)} \\ \text{(O)}}{c:\ \Phi\oplus_\cC(\Psi_\cC\sqcup\Theta_\cC)\ \mid\ \Phi^{e}\oplus_\cC(\Psi_\cC\sqcup\Theta_\cC)}
\end{mathpar}
\noindent
The effect $\Phi\oplus_\cC\Psi$ changes each location of $R_\cC$, read
at the call's entry, within $\Psi$ at its part, and every other location
as $\Phi$ allows. (S), (K) and (O) are proposed static checks of all of
$\cC$'s code: (S)~no \code{await}, \code{yield}, async function or
generator, so the outside call returns inside the call of $\cC$;
(K)~by the flow of the sites each value may come from (Section~\ref{sec:flow}), a value that
crosses, or is stored anywhere but in local storage, reaches no local
site, a crossing value reaches no internal site, and it reaches an
exposed site only at or below a read-only position;
(O)~a store of $v$, an initial field value included, into exposed,
released or outside storage of $R_\cC$, or into storage that the possible
support of an invariant of $G_\cC$ may read, is admitted when all storage
reachable from $v$ is exposed or in $R_\cC$; a store into local storage
outside every invariant's possible support is not constrained, since by (K)
no outside code reaches local storage, a fact that reads through a local
field to unprotected storage is bounded by $\Phi$ there, and only an
invariant's support needs the check, a re-entrant member being able to
redirect a path of it. By (O), a client is
bound only on storage it passed in or received read-only; it is bound
there even through its own writable alias. The checks are of all of
$\cC$'s code: one store that (O) refuses removes the rule for every
outside call inside $\cC$, the walkers' included. Each example below
fails one check; the module is that of Section~\ref{sec:problem} unless
the example says otherwise.

\begin{example}[A suspension]\label{ex:rosuspend}\mbox{}
\begin{center}\small
\begin{tabular}{@{}r@{\hspace{1em}}l@{}}
1 & \code{export async function total(h: H): Promise<number> \{}\\
2 & \code{~~const w = h.widths, n = w.length;}\\
3 & \code{~~await null;}\\
4 & \code{~~let s = 0; for (let i = 0; i < n; i++) s += w[i] as number; return s; \}}\\
\end{tabular}
\end{center}
\emph{Refused by (S)}, at the async function of line 1 and the
\code{await} of line 3. The call returns at the \code{await} while its
frame lives on, so no call of $\cC$ stays active to protect \code{w}: a
caller that runs \code{total(\{ widths: xs \})} and then sets
\code{xs.length = 1} makes it read \code{undefined} at index 1 (observed
under Node 26: the elements read are \code{1, undefined, undefined}).
\end{example}
\begin{example}[A writable value leaked]\label{ex:roleaked}\mbox{}
\begin{center}\small
\begin{tabular}{@{}r@{\hspace{1em}}l@{}}
1 & \code{let last: number[] = [];}\\
2 & \code{export function prepare(n: number): H \{}\\
3 & \code{~~const widths: number[] = []; for (let i = 0; i < n; i++) widths.push(1);}\\
4 & \code{~~last = widths; return \{ widths \}; \}}\\
5 & \code{export function lastWidths(): number[] \{ return last; \}}\\
\end{tabular}
\end{center}
\emph{Refused by (K)}, with \code{walk} of Section~\ref{sec:problem}.
The array of line 3 crosses out read-only in the handle (line 4) and
writable from \code{lastWidths} (line 5). Marked exposed, its crossing at
\code{number[]} on line 5 reaches it at a writable position, which (K)
refuses; marked released, it is outside $R_\cC$, and (O) refuses the
record of line 4 that installs it in protected storage. A visitor that
runs \code{lastWidths().length = 0} makes \code{walk} read
\code{undefined} (observed under Node 26: the sum is \code{NaN}).
\end{example}
\begin{example}[A caller's array installed]\label{ex:roinstall}\mbox{}
\begin{center}\small
\begin{tabular}{@{}r@{\hspace{1em}}l@{}}
1 & \code{const handles: \{ widths: readonly number[] \}[] = [];}\\
2 & \code{export function prepare(n: number): H \{}\\
3 & \code{~~const widths: number[] = []; for (let i = 0; i < n; i++) widths.push(1);}\\
4 & \code{~~const r = \{ widths \}; handles.push(r); return r; \}}\\
5 & \code{export function install(i: number, xs: number[]): void \{}\\
6 & \code{~~const r = handles[i]; if (r !== undefined) r.widths = xs; \}}\\
\end{tabular}
\end{center}
\emph{Refused by (O)}, at the store of line 6. The record is exposed,
so its field is in $R_\cC$, and \code{xs} crossed in at the writable
type \code{number[]}: outside storage that no crossing reached
read-only. Admitted, the store would bind the client on an array it
handed over writable: after \code{install(0, xs)}, a visitor that sets
\code{xs.length = 0} makes \code{walk} read \code{undefined} (observed
under Node 26: the sum is \code{NaN}).
\end{example}
\begin{example}[Released storage coming back]\label{ex:roback}\mbox{}
\begin{center}\small
\begin{tabular}{@{}r@{\hspace{1em}}l@{}}
1 & \code{export const walk = (h: unknown, visit: () => void): number => \{}\\
2 & \code{~~const rs = h as readonly (readonly number[])[]; visit();}\\
3 & \code{~~const row = rs[rs.length - 1]; if (row === undefined) return 0;}\\
4 & \code{~~const n = row.length; let s = 0;}\\
5 & \code{~~for (let i = 0; i < n; i++) \{ visit(); s += row[i] as number; \}}\\
6 & \code{~~return s; \};}\\
\end{tabular}
\end{center}
\emph{Refused by (O)}, with \code{held}, \code{rows} and \code{fill} of
Example~\ref{ex:roreleased} in the same module, at \code{fill}'s store
\code{held.push(row)}. The parameter of line 1 is read at the read-only
type of line 2, so \code{held}, passed back as
\code{walk(rows(), visit)}, is released storage that a crossing in
reaches read-only: it is in $R_\cC$ during the call. A visitor's call of
\code{fill} then stores a row it holds writable into it, which (O)
refuses. A visitor that calls \code{fill(3)} and then empties the row
makes \code{walk} read \code{undefined} (observed under Node 26: the sum
is \code{NaN}).
\end{example}
\begin{example}[A visitor typed to return nothing]\label{ex:rovoid}\mbox{}
\begin{center}\small
\begin{tabular}{@{}r@{\hspace{1em}}l@{}}
1 & \code{export const walk = (h: H, visit: () => void): number => \{}\\
2 & \code{~~const got: unknown = visit(); const r = handles[0];}\\
3 & \code{~~if (r !== undefined \&\& Array.isArray(got)) r.widths = got;}\\
4 & \code{~~const n = h.widths.length; let s = 0;}\\
5 & \code{~~for (let i = 0; i < n; i++) \{ visit(); s += h.widths[i] as number; \}}\\
6 & \code{~~return s; \};}\\
\end{tabular}
\end{center}
\emph{Refused by (O)}, with \code{handles} and \code{prepare} of
Example~\ref{ex:roinstall}, at the store of line 3. A visitor typed
\code{() => void} may return any value, so its result is read at no
read-only type: an array it returns is outside storage that no crossing
reached read-only, and storing it into the handle is refused. A visitor
that returns its own array and empties it at the next visit makes
\code{walk} read \code{undefined} (observed under Node 26: the sum is
\code{NaN}).
\end{example}
\begin{example}[A lookalike slice]\label{ex:roslice}\mbox{}
\begin{center}\small
\begin{tabular}{@{}r@{\hspace{1em}}l@{}}
1 & \code{export type Sliced = \{ slice(): number[] \};}\\
2 & \code{export function install(i: number, xs: Sliced): void \{}\\
3 & \code{~~const r = handles[i]; if (r !== undefined) r.widths = xs.slice(); \}}\\
\end{tabular}
\end{center}
\emph{Refused by (O)}, with \code{handles} and \code{prepare} of
Example~\ref{ex:roinstall}, at the store of line 3. Here \code{xs.slice()}
calls the caller's own method, code outside $\cC$, so its result is a
value entering $\cC$ at the writable type \code{number[]}, not an array
$\cC$ allocates. After \code{install(0, \{ slice() \{ return a; \} \})}, a
visitor that empties \code{a} makes \code{walk} read \code{undefined}
(observed under Node 26: the sum is \code{NaN}). Example~\ref{ex:roguar}
stores the result of the array \code{slice}, which allocates.
\end{example}

\begin{proposition}[Read-only reliance]\label{prop:rorely}
Let $c$ start in a state $s$ inside a call of $\cC$, meeting (M) and
(P2) of Appendix~\ref{app:ro}. Let $F$ hold at $s$ and $S$ support it,
and let each location $l$ of each $(r:\delta)\in S$ meet
$(\Psi_\cC\sqcup\Theta_\cC)(\mathrm{part}(l))\sqsubseteq\delta$ if
$l\in R_\cC(s)$, and $\gamma\sqsubseteq\delta$ for each entry
$(r':\gamma)$ of $\Phi$ with $l\in r'$ otherwise. Then $F$ holds when
$c$ returns (the \emph{bound} route), and so does each invariant of
$G_\cC$ that holds at $s$ and whose support, read at $s$, lies in
$R_\cC$ and allows $\Theta_\cC$, that is,
$\Theta_\cC(\mathrm{part}(l))\sqsubseteq\delta$ for each location $l$ of
each $(r:\delta)$ of the support (the \emph{guarantee} route).
\end{proposition}
\begin{proof}
The rule makes the call an action of $\Phi\oplus_\cC(\Psi_\cC\sqcup\Theta_\cC)$,
which changes each location of $S$ within its $\delta$. The invariants
are (L4).
\end{proof}

\noindent
The bound route serves a fact whose locations no member changes beyond
what the fact allows; the guarantee route a fact a member may rewrite,
since each member's write to the invariant's support is checked to keep
it. Examples~\ref{ex:rocursor}, \ref{ex:roin}, \ref{ex:ro}
and~\ref{ex:rotrim} use the bound route, Example~\ref{ex:roguar} the
guarantee route.

\noindent
Why the rule holds: a step outside $\cC$ reaches no local storage and
the rest of $R_\cC$ only read-only, so by Assumption~\ref{ass:clients}
it changes $R_\cC$ within $\Theta_\cC$; an outside run changes $R_\cC$
otherwise only in members, within $\Psi_\cC$, by induction on the nesting
of calls; and $R_\cC$ only grows inside a call of $\cC$, so the bounds
compose, rigidly when $\Phi$ is rigid and $\Psi_\cC\sqcup\Theta_\cC\sqsubseteq\Phi$
(Appendix~\ref{app:ro})
\status{Lean: \lean{readOnly\_rule}, \lean{readOnly\_frame}, \lean{readOnly\_rely} (the bound route, its premise \lean{Framed} location by location), \lean{readOnly\_guarantee} (the guarantee route, \lean{InvIn}, \lean{SupportAt}), over runs of outside steps meeting \lean{OutsideStep} and member calls meeting \lean{Member} (from \lean{memberOfChecks}). For the calculus's steps at the roots of outside code, \lean{outsideStep\_reach}, \lean{outsideStep\_writes} and \lean{outsideStep\_noForge}, chains following what a closure's code holds (\lean{capture\_reached}), and \lean{rootsValid\_of\_term}. Assumed: the client obligation; that an outside step keeps the marks and the history; that a state's roots are its term's; the members' certified effects; and that the front's description interprets each step of $\cC$'s code, (S), (K) and (O) being run over every body of a component (\lean{ComponentChecks}) and giving a step's checks under that interpretation (\lean{crossOk\_of\_check}, \lean{localStores\_of\_check}, \lean{localBindings\_of\_check}, \lean{admitted\_of\_check}), \lean{frameCheck\_sound} proving the premise's check sound, for mixed supports, given its two inputs: the possible parts cover every location (\lean{flow\_of\_sites}, from a checked flow of sites, Section~\ref{sec:flow}) and every lent location claimed protected has a live history witness (\lean{prot\_of\_witness}). conj.: a run of outside code defined as the calculus's steps}.

\section{The flow of sites}\label{sec:flow}

The check (K) and the premise of Proposition~\ref{prop:rorely} read, for
each value at each node of $\cC$'s code, the allocation sites of the
storage it may reach. A \emph{flow} states them. It is an escape analysis:
its set $O$ holds the sites of the storage outside code may reach. It
reads no type, mark or protocol, so it is stated for any component and any
abstraction of allocations; read-only frames is one of its uses
(Section~\ref{sec:flowro}). An untrusted procedure proposes a flow and a
decision procedure checks it; only the check is trusted.

\paragraph{Sites.} Each allocation of $\cC$'s code, at its occurrence
$p$ in a node expression, a native call's allocation, a closure's and a
rest parameter's array included, has a site $\alpha(p)$ of a finite set
$\Sigma$; all outside storage has the site $\circ$, and
$\Sigma_\circ=\Sigma\cup\{\circ\}$. A location's site is fixed at its
allocation. A \emph{heap of sites} $H:\Sigma_\circ\to\mathcal
P(\Sigma_\circ)$ gives the sites that storage of each site may reference:
through a field, a spine's cells, or a value bound by a closure stored
there; $H(S)=\bigcup_{s\in S}H(s)$. A value $v$ \emph{meets} a set $S$,
$v\models S$, when every location $v$ reaches in one step, through a field,
a cell or what its code holds, has its site in $S$. A heap \emph{meets}
$H$ when the value of each location meets $H$ of the location's site. A
set $S$ is \emph{closed} when $H(S)\subseteq S$; $\overline S$ is the least
closed set containing $S$
\status{Lean: \lean{SiteHeap}, its contents and its bindings apart; \lean{SitesOf}, \lean{Leaves}, \lean{HeapSites}, \lean{Closed}; the site at allocation, \lean{Ghost.alloc}. The Lean's $\alpha$ is constant on each code of $\cC$: one site for its objects and one for its arrays (\lean{Alloc}, \lean{CodeFlow.alloc}, \lean{stepSite})}.

\begin{proposition}[Reach]\label{prop:reach}
In a heap that meets $H$, a value that meets a closed set $S$ reaches only
locations whose sites are in $S$.
\end{proposition}
\begin{proof}
By induction on the chain of reads: the first location is in $S$ by
$v\models S$, and the value stored at a location of site $s\in S$ meets
$H(s)\subseteq S$
\status{Lean: \lean{reach\_sites}}.
\end{proof}

\paragraph{Flows.} A flow $F=(O,H,\tau,\nu,\phi,\rho)$ gives the set
$O\subseteq\Sigma_\circ$ of the sites outside code may reach, a heap of
sites $H$, and for each node $l$ of each code $m$ of $\cC$ the sites
$\tau_l(x)$ of each slot $x$ of its environment, the sites $\nu_l$ of its
expression's value, and, where a slot holds closures of one code of
$\cC$, that code $\phi_l(x)$; $\rho_m$ is the set a return of $m$ meets. A
code is \emph{open} when outside code may call a closure of it: it is a
member's, or a closure of it may meet $O$
\status{Lean: \lean{Flow}, \lean{CodeFlow} (\lean{env}, \lean{val}, \lean{fn}, \lean{vfn}), \lean{FnSig}. The Lean treats every code as open, with $\rho_m=O$}.

\paragraph{Typing by sites.} $\Gamma\vdash e:S$ states that the value of
$e$, once computed, meets $S$, where $\Gamma$ gives the sites of the slots
and variables. At a position of outside code every set lies within $O$.
The rules of the allocating, reading, writing, calling and throwing terms
are
\begin{mathpar}
\inferrule*[left=New]{\Gamma\vdash e_i:S_i \\ S_i\subseteq H(\alpha(p)) \\ \alpha(p)\in S}{\Gamma\vdash\kw{new}^{p}\{\bar f:\bar e\}:S}
\and
\inferrule*[left=Read]{\Gamma\vdash e:S_e \\ H(\overline{S_e})\subseteq S}{\Gamma\vdash e[k]:S}
\and
\inferrule*[left=Write]{\Gamma\vdash e:S_e \\ \Gamma\vdash e':S' \\ S'\subseteq H(s)\ \text{for}\ s\in S_e}{\Gamma\vdash e[k]:=e':S}
\and
\inferrule*[left=Call-out]{\Gamma\vdash f:O \\ \Gamma\vdash e_r:O \\ \Gamma\vdash e_i:O \\ \alpha(p)\in O \\ O\subseteq S}{\Gamma\vdash f(e_r;\bar e)^{p}:S}
\and
\inferrule*[left=Call-in]{\Gamma(f)\text{ names }m \\ \Gamma\vdash e_r:\tau^m(\kw{this}) \\ \Gamma\vdash e_i:\tau^m(x_i) \\ \rho_m\subseteq S}{\Gamma\vdash f(e_r;\bar e)^{p}:S}
\and
\inferrule*[left=Throw]{\Gamma\vdash e:O}{\Gamma\vdash\kw{throw}\;e:S}
\end{mathpar}
with $\tau^m$ the slots of $m$'s entry node, a rest parameter's collecting
the actuals past the parameters in an array of site $\alpha(p)$. A read
meets the closure $\overline{S_e}$, since it may follow prototypes; the
rest spine of \textsc{Call-out} is allocated for a callee that may be
outside code. A variable meets the set of its binder; a test, an
arithmetic operation and a length give a value with no location, which
meets every set; a nested run of $\cC$'s code fits when its environment
meets its node's slots, and of outside code when its environment meets
$O$
\status{Lean: \lean{Fits}, at a value judgment \lean{VJ}: \textsc{New} is \lean{Fits.makeAbsent}, \lean{Fits.makeView}, \lean{Fits.makeBound}, \lean{Fits.makeGraphC}, \lean{Fits.makeGraphOut}, \lean{Fits.arrLit}, \lean{Fits.makeKey}; \textsc{Read} is \lean{Fits.readAt}, \lean{Fits.indexAt}, \lean{Fits.elemAt}, \lean{Fits.spread}, \lean{Fits.merge}; \textsc{Write} is \lean{Fits.writeAt}, \lean{Fits.setElem}; \textsc{Call-out} is \lean{Fits.invoke}, \lean{Fits.construct}; \textsc{Call-in} is \lean{Fits.invokeC}, whose actuals are values or variables, at the code's array site, with $\rho_m=O$; \textsc{Throw} is \lean{Fits.throw}; runs \lean{Fits.frameC}, \lean{Fits.frameOut}, \lean{Fits.gotoC}, \lean{Fits.gotoOut}}.

\begin{definition}[A flow of $\cC$]\label{def:flow}
$F$ is a flow of $\cC$ when
(F1) each node's expression meets its value set, $\tau_l\vdash e_l:\nu_l$;
(F2) each edge from $l$ to $l'$ carries its values: each slot of $l'$ holds
the sites of the slot or the value it is filled from, and the slot a handler
binds holds $O$;
(F3) a returning node's value meets $\rho_m$, and $\rho_m\subseteq O$ for
an open code;
(F4) $\circ\in O$, $O$ is closed, and storage of a site of $O$ may hold
every site of $O$;
(F5) the entry slots of an open code hold $O$
\status{Lean: \lean{CodeOk}, \lean{ContOk}, \lean{HandlerOk}, \lean{Flow.Facts}; with (F5) for every slot of every node (\lean{Flow.Facts.envO})}.
\end{definition}

\begin{theorem}[Flow]\label{thm:flow}
Let $F$ be a flow of $\cC$, and let each call of $\cC$ be entered with
well-formed values. At every state reached in that call, the heap meets $H$
and every value held by outside code meets $O$; so, by
Proposition~\ref{prop:reach}, every location outside code reaches has its
site in $O$.
\end{theorem}
\begin{proof}
Every step keeps the term typed and the heap meeting $H$, the locations it
allocates at the site of the allocating occurrence: by induction on the
step, a read giving what $H$ allows at a site the reach may have, a write
and an allocation storing what $H$ allows, a call entering its callee's
slots, a return giving $O$. A step of outside code changes only what its
roots reach, stores there only what they reach or what it allocates, and
allocates outside storage, so (F4) keeps both. At an entry every location
is outside storage, of site $\circ\in O$
\status{Lean: \lean{fits\_step}, \lean{sites\_outside}, \lean{ClientFrame.fitsAt\_reached}, \lean{fits\_roots}, \lean{ClientFrame.sitesSound\_of\_fits}; the entry, \lean{ClientFrame.entryFits} from \lean{EntryShape}, well-formed entries assumed. Proved where every code is open and $\alpha$ is constant on each code; conj.\ for the statement as given}.
\end{proof}

\paragraph{Proposal and check.} A flow is proposed by iterating the rules
to a post-fixed point over the nodes of $\cC$'s codes: the least set of
each node's expression at its slots, its stores added to $H$, what it
gives outside code added to $O$, and the actuals of a call in added to
the callee's entry slots, until nothing changes. Nothing is proved of the
proposal. The check decides Definition~\ref{def:flow} of a finite flow: it
computes each node expression's least set from its slots and compares it
with $\nu_l$; it compares each edge's slots up to a width past which both
ends are constant; and it decides (F4) and (F5).

\begin{theorem}[Check]\label{thm:check}
A finite flow the check accepts is a flow of $\cC$.
\end{theorem}
\noindent{\sloppy\status{Lean: \lean{FFlow.check\_sound}, from \lean{syn\_sound} (the least set fits, and so does every larger one), \lean{codeChkB\_sound}, \lean{carriesBaseB\_sound} and \lean{codeOk\_of\_chk}; the proposal is \lean{solveFlow}, untrusted. The check is not complete: it may refuse a flow of $\cC$.}\par}

\smallskip\noindent
The flow and its check have a general form: where every step from a
state the abstraction describes at a node is covered by a constraint into
a node whose transfer describes the next state, a proposal that meets each
constraint, checked by one comparison each, describes every state of an
execution \status{Lean: \lean{checkPost\_sound}, \lean{realizes\_along}, for any abstract domain and constraints; Theorem~\ref{thm:flow} is proved directly, by subject reduction, not as their instance}.

\paragraph{Other uses.} Two values that meet disjoint closed sets reach no
common location (Proposition~\ref{prop:reach}). A rule of \emph{Tolerant
frames} that needs two paths apart could read it from the slots of a node
(conj.: it needs the slots' sets closed under $H$ and a support's regions
read as for the frame below). The check (O) could read the sites of a
stored value from the flow, admitting a store whose value meets a closed
set outside $O$ (conj.: it needs (O)'s classification of the stores
stated over sites).

\subsection{Read-only frames from the flow}\label{sec:flowro}

\begin{definition}[The marks of a flow]
A site of $O$ is released, and every other site of $\Sigma$ local.
\end{definition}
\noindent
By Proposition~\ref{prop:kflow}, the storage of a site outside $O$ is held only by $\cC$'s frames and
by local storage, and lies in $R_\cC$ at every state; the storage of a
site of $O$ lies in $R_\cC$ while a crossing in has reached it read-only.
The flow tells nothing of types, so its marks are local and released only:
a site that outside code reaches only read-only, which (K) admits as
exposed, is released here.

\begin{proposition}[The check (K) from a flow]\label{prop:kflow}
With the marks of a flow $F$ of $\cC$, where no closure of outside code
or bound function stored in $\cC$'s storage binds a value of a local site,
every step of $\cC$'s code from a reached state meets (K).
\end{proposition}
\begin{proof}
A value that crosses out at a step is held by outside code after it, or
left in storage that outside code reaches, of a site of $O$, whose
contents $O$ holds (F4); by Theorem~\ref{thm:flow} it meets $O$, which
holds released sites and $\circ$ only. Storage of a released site
references sites of $O$ only, since $O$ is closed, so it references no
local site
\status{Lean: \lean{ClientFrame.chkSound\_of\_flow}, through \lean{escapeAt\_of\_sites}, \lean{crossOk\_of\_sites}, \lean{localStores\_of\_sites}, \lean{localBindings\_of\_sites}, at the actual crossings of a step (\lean{actualXs}, \lean{gives\_actual}, \lean{stores\_actual}). The Lean takes the marks as data and decides that they meet $O$ (\lean{FFlow.marksB\_sound}: \lean{Lendable}, \lean{LocalHeld})}.
\end{proof}

\noindent
At a call of outside code at node $l$, each region $r$ of a support has
the sites $\Sigma_l(r)$: a field $x.f$ the sites $\tau_l(x)$, a path
$x.\pi$ the sites that $H$ gives one field after another, a spine those of
its array, its cells $H$ of them. Each location of $r$ has its site in
$\Sigma_l(r)$, by Theorem~\ref{thm:flow}. The check of the premise of
Proposition~\ref{prop:rorely} at $l$ is, for each $(r:\delta)$ of the
support and each $s\in\Sigma_l(r)$,
\[
\begin{cases}
(\Psi_\cC\sqcup\Theta_\cC)(\mathrm{part})\sqsubseteq\delta & \text{if } s\notin O, \text{ or } r \text{ is read as protected},\\
\gamma\sqsubseteq\delta\ \text{for each}\ (r':\gamma)\in\Phi\ \text{not apart from}\ r & \text{otherwise,}
\end{cases}
\]
at the part of $r$'s locations of site $s$; a region read as protected
needs, for each of its locations of a site of $O$, a crossing in of the
active call that reached it read-only.

\begin{proposition}[The frame from a flow]\label{prop:frameflow}
At a call of outside code from a frame of $\cC$'s code at node $l$, the
check above gives the premise of Proposition~\ref{prop:rorely}.
\end{proposition}
\noindent{\sloppy\status{Lean: \lean{regionSites}, \lean{regionSites\_sound}, \lean{frameCheckAt}, \lean{framed\_of\_sites} (through \lean{frameCheck\_sound}, \lean{flow\_of\_sites} and \lean{prot\_of\_witness}), \lean{framed\_of\_frame}, at the innermost frame of a reached state \lean{ClientFrame.framed\_at\_innermost\_keys}; the crossing in as a live history witness, \lean{Witness}, \lean{entryStart\_witness}, \lean{witness\_along}. The Lean reads every field of a site as one part (\lean{noKeys}), and gives no sites to a region of a class, of every location or of a value's fields by enumeration, which the check then refuses}.\par}

\paragraph{The examples by the flow.}
\begin{itemize}[leftmargin=1.5em,itemsep=2pt]
\item Example~\ref{ex:rocursor}: only \code{walk}'s frame holds \code{c},
and \code{visit(c.i)}, a call out, passes a number, which meets every set.
The cursor's site is outside $O$, local, and $c.i$ lies in $R_\cC$. The
array \code{xs} fills an entry slot, which holds $O$ by (F5); its spine,
outside storage, is not read as protected, and $\Phi$ bounds it (measured
2026-11-19 by the flow probe: the check accepts the proposal for
\code{walk}, the cursor's site outside $O$).
\item Example~\ref{ex:roreleased}: \code{rows} and \code{fill} return
\code{held} and \code{row}, so by (F3) both sites are in $O$, released.
\item Section~\ref{sec:problem}'s module and Examples~\ref{ex:roin}
and~\ref{ex:ro}: \code{prepare} returns its record and its array, so their
sites are in $O$, released. The spine of \code{w} is read as protected,
the crossing of \code{h} at \code{H} its witness, whoever allocated it.
\item Example~\ref{ex:roleaked}: \code{lastWidths} returns the array, so
its site is in $O$, and the example's argument for a released array, by
(O), applies.
\item Examples~\ref{ex:roinstall} to~\ref{ex:roslice} are refused by (O),
which the flow does not decide; it gives the sites they read. The
visitor's result in Example~\ref{ex:rovoid} and the result of the caller's
\code{slice} in Example~\ref{ex:roslice} are results of calls out, of
$O$.
\item Example~\ref{ex:roguar}: the array \code{slice} allocates at a site
of $\cC$, and the copy is stored into a field of \code{h}, of a site of
$O$, so by \textsc{Write} and (F4) its site is in $O$, released. (O)
admits the store at the mark exposed only, which the flow does not
confirm.
\item Examples~\ref{ex:rosuspend}, \ref{ex:rotrim} and~\ref{ex:rocast}
turn on (S), the protocol and the client assumption, which read no flow.
\end{itemize}
Two more examples show the two calls.

\begin{example}[A cursor passed to a helper]\label{ex:rohelper}\mbox{}
\begin{center}\small
\begin{tabular}{@{}r@{\hspace{1em}}l@{}}
1 & \code{type Visitor = (w: number) => Modifies<\{ values: "any"; spines: "grow" \}, void>;}\\
2 & \code{function next(c: \{ i: number \}): void \{ c.i++; \}}\\
3 & \code{export function walk(xs: number[], visit: Visitor): number \{}\\
4 & \code{~~const c = \{ i: 0 \}; let s = 0;}\\
5 & \code{~~while (c.i < xs.length) \{ visit(c.i); s += xs[c.i]!; next(c); \}}\\
6 & \code{~~return s; \}}\\
\end{tabular}
\end{center}
\emph{Accepted}, as Example~\ref{ex:rocursor}. The slot of \code{walk} that
holds \code{next} names its code, so \code{next(c)} on line 5 is a call
in: \code{c} meets \code{next}'s entry slot, and no rule puts the
cursor's site in $O$. \code{next} is not open, so its slots need not hold
$O$. Pretext's \code{walkPreparedLinesRaw} has this shape (measured
2026-11-19 by the flow probe: the check accepts its proposal, its cursor's
site outside $O$).
\end{example}

\begin{example}[A cursor given away]\label{ex:rogiven}\mbox{}
\begin{center}\small
\begin{tabular}{@{}r@{\hspace{1em}}l@{}}
1 & \code{type Cursor = \{ i: number \};}\\
2 & \code{type Visitor = (c: Cursor) => Modifies<\{ values: "any"; spines: "grow" \}, void>;}\\
3 & \code{export function walk(xs: number[], visit: Visitor): number \{}\\
4 & \code{~~const c: Cursor = \{ i: 0 \}; let s = 0;}\\
5 & \code{~~while (c.i < xs.length) \{ visit(c); s += xs[c.i]!; c.i++; \}}\\
6 & \code{~~return s; \}}\\
\end{tabular}
\end{center}
\emph{Refused by the premise} of Proposition~\ref{prop:rorely}.
\code{visit(c)} on line 5 is a call out, so \code{c} meets $O$ and its site
is in $O$, released; no crossing in reaches it, so $c.i$ is not read as
protected, and $\Phi$ allows $\any\not\sqsubseteq\keep$ there. A visitor
that sets \code{c.i = 5} makes \code{walk} read \code{undefined} (observed
under Node 26: the sum is \code{NaN}).
\end{example}
\section{Examples and limits}\label{sec:examples}

\begin{example}[The module above]\label{ex:ro}
\code{visit()} has the effect $\top\oplus_\cC\keep$. The support
$\{w.\kw{spine}:\grow\}$ of $i<\len w$ is in $R_\cC$: exposed when
\code{prepare} allocated it, under the marks of Section~\ref{sec:concepts},
and protected by the crossing of \code{h} whoever allocated it, and $\keep\sqsubseteq\grow$: $i<\len w$ survives,
under Assumption~\ref{ass:clients} with $\Theta_\cC=\keep$.
\end{example}

\begin{example}[A member that shortens]\label{ex:rotrim}\mbox{}
\begin{center}\small
\begin{tabular}{@{}r@{\hspace{1em}}l@{}}
1 & \code{export function trim(h: H): void \{ (h.widths as number[]).pop(); \}}\\
\end{tabular}
\end{center}
\emph{Refused by the protocol.} \code{trim} is a member, and line 1
shortens exposed storage, so a protocol that bounds the members has
$\any$ at that spine; $\any\not\sqsubseteq\grow$ fails the premise of
Proposition~\ref{prop:rorely} for $i<\len w$ in \code{walk}. A visitor
that calls \code{trim(h)} twice makes \code{walk} read \code{undefined}
(observed under Node 26: the elements read are
\code{1, undefined, undefined}).
\end{example}

\begin{example}[A client that casts]\label{ex:rocast}\mbox{}
\begin{center}\small
\begin{tabular}{@{}r@{\hspace{1em}}l@{}}
1 & \code{const h = prepare(3);}\\
2 & \code{walk(h, () => \{ (h.widths as number[]).pop(); \});}\\
\end{tabular}
\end{center}
\emph{Accepted, under an assumption this client breaks.} The derivation
of Example~\ref{ex:ro} holds, and the client's pop on line 2 changes
exposed storage in $R_\cC$ beyond $\Theta_\cC=\keep$, which
Assumption~\ref{ass:clients} excludes. The walker then reads
\code{undefined} (observed under Node 26: the elements read are
\code{1, 1, undefined}); the fault is the client's.
\end{example}

\begin{example}[A guarantee]\label{ex:roguar}
Give \code{H} a second array and an invariant, and a member that installs
a copy:
\begin{center}\small
\begin{tabular}{@{}l@{}}
\code{/** @invariant this.widths.length === this.flags.length */}\\
\code{type H = \{ readonly widths: readonly number[]; readonly flags: readonly number[] \};}\\
\code{export function install(h: H, xs: number[]): void \{}\\
\code{~~if (xs.length === h.flags.length) (h as \{ widths: number[] \}).widths = xs.slice(); \}}
\end{tabular}
\end{center}
A walker that reads \code{h.widths[i]!} after each visit, with
$i<\len{h.\mathit{flags}}$, needs the invariant $I$ after the visit. The
bound route fails: \code{install} replaces the array at the key
\code{widths}, so $\Psi_\cC$ is $\any$ there. The guarantee route holds.
The copy is allocated by $\cC$ and reached by outside code only through
\code{H}, so its site is exposed, and (O) admits the store; by the
contract of the array \code{slice} (\emph{Tolerant frames}, Section~6,
``Library functions''), the copy has the length of \code{xs},
which the test equates with $\len{h.\mathit{flags}}$, so \code{install}
keeps $I$ and $I\in G_\cC$. Read at the visit's entry, $I$'s support (the
two fields at $\grow$ and their two spines at $\keep$) lies in $R_\cC$,
and $\keep$ allows $\Theta_\cC=\keep$. So $I$ holds after the visit, also
when the visitor calls \code{install} and pops its own \code{xs}, which
is outside the support: the walker reads the copy (observed under Node 26).
\end{example}

\noindent
A fact relies on one component's rule; combining two components' rules
is not stated. The model's objects carry no integrity state (frozen,
sealed, non-extensible), and no protocol states its change.

\section{Related work}\label{sec:related}

Each work below is cited for what it establishes, in the version its
bibliography entry names; a correspondence with this note is marked
cited when it compares stated facts of both, and conj.\ when it goes
beyond them.

\paragraph{Hidden invariants of a module.} Banerjee and Naumann~\cite{bn}
state in a module's interface a \emph{dynamic boundary}, a read effect
over ordinary and ghost state that frames the module's hidden invariant,
verify every client command to write none of its locations, and prove a
second-order frame rule sound, a variant for reentrant callbacks
included, for first-order sequential programs. $R_\cC$ is a dynamic
boundary designated through ghost state, Assumption~\ref{ass:clients}
the client's side of it, and \textsc{Read-only frame} the second-order
frame rule (cited). Their clients are all verified and write nothing in
the boundary; here unanalysed clients are assumed, callbacks change
$R_\cC$ within $\Psi_\cC$ by calling members, and a callback passed as a
value, outside their language, is the case at hand (cited). The
hypothetical frame rule of O'Hearn, Yang and Reynolds~\cite{oyr} hides a
module's resource invariant from clients that change it only through the
module's procedures, with the higher-order counterpart of Birkedal,
Torp-Smith and Yang~\cite{bty}; separation makes the hidden storage
unreachable to the client, while here outside code holds references to
exposed storage and reads it (cited). Pottier's anti-frame
rule~\cite{antiframe} hides local state in direct style by linear
capabilities that a type system tracks for all code (cited, from its
abstract).

\paragraph{Object and class invariants.} The Boogie
methodology~\cite{boogieinv} records which invariants hold in auxiliary
fields changed by \code{pack} and \code{unpack}; its callers establish
invariant preconditions, as a handle's invariant here is a contract of
function boundaries (cited). Leino and M\"uller treat invariants over
dynamic ownership contexts~\cite{lmdynamic} and static class invariants
broken only inside a non-reentrant \code{expose}~\cite{lmstatic};
Drossopoulou et al.~\cite{dfms} frame visible-state techniques, and
Summers and Drossopoulou~\cite{considerate} let a method declare the
invariants it may find broken. Each admits an invariant broken between
visible states; here a member's write to the support of an invariant of
$G_\cC$ is checked to keep it, so a callback never finds one broken
(cited). Leino and Nelson~\cite{leinonelson} define modular soundness and
prove it for static dependencies (cited); Assumption~\ref{ass:clients}
is the condition under which a component's checked facts stay valid as
more code is linked around it (conj.).

\paragraph{Ownership and read-only references.} Universes~\cite{universes}
enforce owner-as-modifier: a reference not through an object's owner only
reads; the survey of Clarke et al.~\cite{ownsurvey} places it among the
relaxations of owners-as-dominators. Assumption~\ref{ass:clients} is an
owner-as-modifier condition at a component's granularity (cited), assumed
of unanalysed code rather than typed. Javari~\cite{javari} gives
transitive reference immutability, which matches the transitivity of a
read-only position here but bounds only the reference, where the
assumption bounds the storage, through any alias (cited).

\paragraph{Unknown code and scoped protection.} Swasey, Garg and
Dreyer~\cite{ocpl} prove robust safety: a program that passes only low
values to arbitrary code keeps its invariants, a read-only wrapper among
their examples. Here outside code holds direct references to protected
storage, and the guarantee holds only for outside code meeting
Assumption~\ref{ass:clients} (cited). Drossopoulou et al.~\cite{external}
call an object protected when no external object reachable from the
current frame accesses it, and a scoped invariant holds until the call
active in a state returns; the history of Appendix~\ref{app:ro}, under
which storage protected in a call stays protected until the outermost
call returns, is a scoped future in that sense (cited). Implicit dynamic
frames~\cite{idf} infer a frame from access assertions, as $\Psi_\cC$ is
the members' checked effects, but make a callee's frame explicit at each
call, where an unknown callback here has the effect $\top$ (cited). A
history invariant~\cite{history} bounds how a subject's data may evolve,
as $\Theta_\cC$ bounds what clients change; it is checked of every
update, where $\Theta_\cC$ is assumed of unanalysed code (cited).


\appendix
\section{The proof of Read-only frame}\label{app:ro}

This appendix states the steps of the argument for \textsc{Read-only
frame} (Section~\ref{sec:rule}), for one component $\cC$.

\paragraph{History.} A state carries three ghost components: the set $E$
of exposed locations that a crossing has reached read-only; a stack
$A$ of sets of locations, one for each active call of $\cC$; and a set $B$
of the released and outside locations protected while no call of $\cC$ is
active. At a call of an entry, the call's set, $\bigcup A\cup B$, is pushed
first, and then its arguments cross. A crossing of $v$ at $T$ adds to $E$
the exposed locations reached read-only from $v$, and, when it is
\emph{in}, adds the released and outside locations reached read-only
from $v$ to every set of $A$. Each exit, normal or exceptional, pops the
top set; when the stack becomes empty, $B$ becomes the popped set's
locations that $\cC$ holds. A step while no call of $\cC$ is active
removes from $B$ each location $\cC$ no longer holds. Every set of $A$
is then $\bigcup A\cup B$, so the stack is one set and a depth, and
\[
l\in R_\cC(s)\iff\mu(l)\in\{\kw{local},\kw{exposed}\}\ \vee\ \bigl(\mu(l)\in\{\kw{released},\kw{out}\}\wedge l\in{\textstyle\bigcup}A\cup B\bigr)
\]
with $\mu(l)$ the mark of $l$'s site
\status{Lean: \lean{Ghost.cross}, \lean{HStep}, \lean{HStep.uniform}, \lean{exit\_unprotects}, \lean{Ghost.InR}; over an execution, \lean{Ghost.next}, the removal from $B$ at a step outside every call \lean{outsideNext}; holding is \lean{Held}. The Lean's history has a further transition, \lean{HStep.hold}, adding any set to every set of $A$, which the rule does not use}.

\paragraph{Invariants of states.} A state is \emph{ok} for $\cC$ when a
call of $\cC$ is active and:
\begin{description}[leftmargin=1.5em,style=nextline]
\item[(M) Marks] every location's mark is its site's, fixed at
allocation; a local or internal location is reachable from no frame of
code outside $\cC$; storage that references local storage is local; and
no value that a closure of outside code, or a bound function, held in
storage binds (a capture, a saved receiver or argument) reaches local
storage;
\item[(P2) Confinement] every exposed location reachable from a frame of
code outside $\cC$ is in $E$.
\end{description}
Reachability follows fields, spines, cells and what a closure's code
holds. A closure of $\cC$ binds what it captures, local storage among it
\status{Lean: \lean{MarksOk}, \lean{LocalClosed}, \lean{BindingsClosed}, \lean{Confined}; kept by outside steps (\lean{outside\_localClosed}, \lean{outside\_bindingsClosed}) and by checked steps of $\cC$ (\lean{pkg\_step\_ok}); \lean{binding\_no\_local}, no chain from such a binding reaches local storage}.

\paragraph{Lemmas.} For ok states and outside steps that meet
Assumption~\ref{ass:clients}:
\begin{description}[leftmargin=1.5em,style=nextline]
\item[(L1) Outside step] a step $s\to s'$ outside $\cC$ changes each
$l\in R_\cC(s)$ within $\Theta_\cC$ at its part, and $s'$ meets (M) and
(P2);
\item[(L2) Outside run] a run $s\to^{*}s'$ outside $\cC$ inside an active
call of $\cC$ keeps each $l\in R_\cC(s)$ in the region of each of its
states, changes it within $\Psi_\cC\sqcup\Theta_\cC$ at its part, and
ends in an ok state; by induction on the nesting of calls, a member's
call by its bound;
\item[(L3) Rigidity] for $\Phi$ rigid and $\Psi_\cC\sqcup\Theta_\cC\sqsubseteq\Phi$,
the relation of two states by an action of
$\Phi\oplus_\cC(\Psi_\cC\sqcup\Theta_\cC)$ with a growing history is
transitive, and $R_\cC(s)\subseteq R_\cC(s')$ for related states within a
call, since the region reads the marks, fixed at allocation, and the
history, which only grows inside a call;
\item[(L4) Invariants] an invariant $I\in G_\cC$ that holds at $s$, whose
support lies in $R_\cC$ and allows $\Theta_\cC$, holds after a run as in
(L2): each member's writes and exits keep it, and each outside step
changes its support within $\Theta_\cC$. A walker then reads $I$'s
consequences after the call as at its entry;
\item[(O) Ownership] after a store that (O) admits, the locations
reachable from the stored value are in $R_\cC$ at every later state of
the call.
\end{description}
\noindent{\sloppy\status{Lean: (L1) \lean{outside\_step}; (L2) \lean{Run.bound}; (L3) \lean{HSteps.grows}, \lean{InR.mono}, \lean{Leads.trans}, \lean{Leads.trans\_top}, \lean{Change.Holds.trans}; (L4) \lean{Run.keeps} with \lean{OutsideKeeps}, from \lean{SupportInR} or \lean{SupportAt}, and \lean{InvIn.of\_step} at a member's write; (O) \lean{ownership}, \lean{Admitted}. conj.: a run of outside code defined as the calculus's steps, as for the rule}.\par}

\paragraph{Nested components.} The components of a program form a
laminar family: any two are nested or disjoint. A step outside $\cC$ is
outside each component nested in $\cC$. A site of a component nested in
$\cC$ has a mark for each, its mark for $\cC$ at most its mark for the
nested one in the order
$\kw{local}\sqsubseteq\kw{internal}\sqsubseteq\kw{exposed}\sqsubseteq\kw{released}$,
so the storage of a nested component is never outside storage of $\cC$
\status{Lean: \lean{Family}, \lean{Family.outside\_mono}, \lean{Family.runs\_laminar}, \lean{Family.nested\_origin}, \lean{Family.nested\_local\_inR}, \lean{Family.nested\_exposed}, \lean{MarksConsistent}}.

\paragraph{Client obligations other than keep.} An API may state a
$\Theta_\cC$ other than $\keep$: $\extend$ on a spine lets clients append
to a component's array, $\kw{scalar}$ at a key lets them set a field
that holds no object. The lemmas hold for every $\Theta_\cC$; (L3) asks
of it that $\Psi_\cC\sqcup\Theta_\cC\sqsubseteq\Phi$, and (L4) that the
invariant's support allows it.



\begin{thebibliography}{9}\small\itemsep1pt
\bibitem{tf} TScert working note. \emph{Tolerant frames: a modal account of fact preservation without ownership}. 11 November 2026, revised 17 November.
\bibitem{bn} A.~Banerjee, D.~Naumann. Local reasoning for global invariants, part II: dynamic boundaries. J.~ACM 60(3), 2013.
\bibitem{boogieinv} M.~Barnett, R.~DeLine, M.~F\"ahndrich, K.~R.~M.~Leino, W.~Schulte. Verification of object-oriented programs with invariants. J.~Object Technology 3(6), 2004.
\bibitem{bty} L.~Birkedal, N.~Torp-Smith, H.~Yang. Semantics of separation-logic typing and higher-order frame rules. LICS 2005.
\bibitem{ownsurvey} D.~Clarke, J.~\"Ostlund, I.~Sergey, T.~Wrigstad. Ownership types: a survey. In \emph{Aliasing in Object-Oriented Programming}, LNCS 7850, 2013.
\bibitem{universes} W.~Dietl, P.~M\"uller. Universes: lightweight ownership for JML. J.~Object Technology 4(8), 2005.
\bibitem{dfms} S.~Drossopoulou, A.~Francalanza, P.~M\"uller, A.~J.~Summers. A unified framework for verification techniques for object invariants. ECOOP 2008.
\bibitem{external} S.~Drossopoulou, J.~Mackay, S.~Eisenbach, J.~Noble. Reasoning about external calls. PACMPL 9(OOPSLA2), 2025.
\bibitem{lmdynamic} K.~R.~M.~Leino, P.~M\"uller. Object invariants in dynamic contexts. ECOOP 2004.
\bibitem{lmstatic} K.~R.~M.~Leino, P.~M\"uller. Modular verification of static class invariants. FM 2005.
\bibitem{leinonelson} K.~R.~M.~Leino, G.~Nelson. Data abstraction and information hiding. ACM TOPLAS 24(5), 2002.
\bibitem{history} K.~R.~M.~Leino, W.~Schulte. Using history invariants to verify observers. ESOP 2007.
\bibitem{oyr} P.~O'Hearn, H.~Yang, J.~Reynolds. Separation and information hiding. ACM TOPLAS 31(3), 2009.
\bibitem{antiframe} F.~Pottier. Hiding local state in direct style: a higher-order anti-frame rule. LICS 2008.
\bibitem{idf} J.~Smans, B.~Jacobs, F.~Piessens. Implicit dynamic frames: combining dynamic frames and separation logic. ECOOP 2009.
\bibitem{considerate} A.~J.~Summers, S.~Drossopoulou. Considerate reasoning and the composite design pattern. VMCAI 2010.
\bibitem{ocpl} D.~Swasey, D.~Garg, D.~Dreyer. Robust and compositional verification of object capability patterns. PACMPL 1(OOPSLA), 2017.
\bibitem{javari} M.~S.~Tschantz, M.~D.~Ernst. Javari: adding reference immutability to Java. OOPSLA 2005.
\end{thebibliography}

\end{document}
